% !TeX root = ../初等数论笔记.tex
% 由主文档通过 \input 载入；请编译主文档。
\setcounter{chapter}{0} % 保留原书章号。
\chapter{绪论与整数证明方法}
\label{ch:introduction}

初等数论中的计算对象很熟悉，论证方式却常与代数运算题不同。检验一百个例子仍不能证明“对所有整数都成立”，而一个反例就足以否定全称命题。本章从整数的离散性出发，说明归纳、反证、构造和抽屉原理各适合解决什么问题，为后面的证明建立共同语言。

\section{从三个问题认识数论}
\label{sec:intro-questions}

每隔 $6$ 天与每隔 $8$ 天举行的两项活动，若今天同时举行，则下一次同时举行在 $24$ 天后。这是公倍数问题。若两项活动的起始时间不同，就需要同时满足两个余数条件，这将引出同余方程。

整数 $65$ 有表示 $65=1^2+8^2=4^2+7^2$，而 $21$ 不能表示为两个整数的平方和。分数 $1/7=0.\overline{142857}$ 的循环节长度是 $6$。这些现象分别引出平方和判据和乘法阶。数论不仅寻找一个例子，更要回答：何时存在，是否唯一，怎样求出全部对象，以及计算过程为什么正确。

\section{整数、逻辑语言与符号约定}

记 $\mathbb Z$、$\mathbb Q$、$\mathbb R$ 分别为整数、有理数、实数集，$\mathbb Z_{>0}$ 和 $\mathbb Z_{\ge0}$ 分别为正整数集和非负整数集。后文还使用复数集 $\mathbb C$。字母的范围由所在命题规定，通常用 $p,q$ 表示素数，但约定不代替假设。有限集合 $A$ 的元素个数记为 $\#A$。

命题 $P\Rightarrow Q$ 表示 $P$ 是 $Q$ 的充分条件，$Q$ 是 $P$ 的必要条件；$P\Leftrightarrow Q$ 要求两个方向都成立。逆命题 $Q\Rightarrow P$ 一般不成立，逆否命题 $\neg Q\Rightarrow\neg P$ 则与原命题等价。例如，被 $4$ 整除的整数一定是偶数，但偶数 $6$ 不能被 $4$ 整除。

量词顺序不可随意交换。命题“对每个整数 $a$，存在整数 $b$ 使 $b>a$”成立，但不存在一个整数 $b$ 大于所有整数。否定全称命题只需一个反例；否定存在命题则须排除全部可能。

例如，“每个偶数都是两个奇数的和”可写为：对任意 $n\in\mathbb Z$，存在奇整数 $u,v$，使 $2n=u+v$。取 $u=1,v=2n-1$ 就完成了构造；若额外要求 $u,v$ 为正数，就必须另加 $n\ge1$，因为负偶数不可能是两个正数的和。看似细小的取值范围会改变命题本身，所以应先读清对象，再开始运算。

后文的符号 $\sum$ 和 $\prod$ 分别表示求和与求积。例如 $\sum_{j=1}^r a_j=a_1+\cdots+a_r$，$\prod_{j=1}^r a_j=a_1\cdots a_r$；空和约定为 $0$，空积约定为 $1$。这些约定使没有素因子或没有待加项的边界情形也能使用统一公式。

\section{良序原理与数学归纳法}

\begin{axiom}[良序原理]
每个非空的非负整数子集都有最小元。
\end{axiom}

\emph{良序原理（well-ordering principle）}是本书对整数所使用的基本性质。它既用于选择最小对象，也用于证明下降过程必定终止。

\begin{theorem}[强归纳法（strong induction）]
设 $n_0\in\mathbb Z_{\ge0}$。若命题 $P(n_0)$ 成立，且对每个 $n>n_0$，由所有 $n_0\le k<n$ 的 $P(k)$ 成立可推出 $P(n)$，则 $P(n)$ 对所有 $n\ge n_0$ 成立。
\end{theorem}
\begin{proof}
若有反例，取最小反例 $m$。由初始情形知 $m>n_0$，由最小性知 $P(k)$ 对所有 $n_0\le k<m$ 成立；归纳条件于是给出 $P(m)$，矛盾。
\end{proof}

普通归纳法只使用前一个命题 $P(n-1)$，是上述方法的特例。最小反例法把同一论证反过来组织：假定反例存在，再从最小反例构造更小的反例。任何严格递减的非负整数序列都只能有有限项，否则其值集有最小元，而下一项更小，发生矛盾。

归纳假设是一个临时前提，用于证明从已知情形走向下一情形的逻辑蕴含；它不等于预先假定整个结论成立。若构造 $n$ 的对象时只需使用 $n-1$ 的对象，普通归纳通常足够；若需要把 $n$ 拆成两个都小于 $n$、但未必等于 $n-1$ 的数，强归纳往往更自然。第\ref{ch:primes}章的素因子分解正属于后一种情形。

\begin{example}
证明对每个 $n\ge1$，$1+3+\cdots+(2n-1)=n^2$。
\end{example}
\begin{proof}
$n=1$ 时成立。若前 $n$ 项之和为 $n^2$，则前 $n+1$ 项之和为 $n^2+(2n+1)=(n+1)^2$，归纳完成。
\end{proof}

\section{反证、构造与抽屉原理}

\begin{example}
$\sqrt2$ 不是有理数。
\end{example}
\begin{proof}
若 $\sqrt2=a/b$，可在所有 $a,b>0$ 的整数表示中选取分母 $b$ 最小的一组。由 $a^2=2b^2$ 知 $a$ 为偶数：奇数 $2k+1$ 的平方仍为奇数。写 $a=2c$，得 $b^2=2c^2$，故 $b=2d$。于是 $\sqrt2=c/d$ 且 $0<d<b$，与选择矛盾。
\end{proof}

该证明同时说明了反证与递降的作用。构造性证明则必须实际给出对象，并验证它满足条件；仅说“可以找到”不构成存在性证明。

\begin{proposition}[抽屉原理]
若 $N>k\ge1$，把 $N$ 个对象放入 $k$ 个盒子，则至少一个盒子含两个对象。
\end{proposition}
\begin{proof}
若每盒至多一个，对象总数至多为 $k$，与 $N>k$ 矛盾。
\end{proof}

\emph{抽屉原理（pigeonhole principle）}在数论中常以余数为盒子。两个对象落在同一盒子，意味着它们的差被某个整数整除；这一思想将在同余与平方和问题中反复出现。

\begin{example}[从部分和中寻找整除关系]
给定任意 $n$ 个整数 $a_1,\ldots,a_n$，其中 $n\ge1$。证明存在一段连续项，其和被 $n$ 整除。
\end{example}
\begin{proof}
考虑 $n+1$ 个部分和 $S_0=0$、$S_j=a_1+\cdots+a_j$（$1\le j\le n$）。整数除以 $n$ 只有 $n$ 种余数，因此有 $0\le i<j\le n$ 使 $S_i,S_j$ 余数相同。于是 $S_j-S_i=a_{i+1}+\cdots+a_j$ 被 $n$ 整除。这里加入 $S_0$，使从第一项开始的连续段也包含在统一论证中。
\end{proof}

\section{从观察到证明}

发现规律与证明规律是两个相互配合的步骤。例如，前几个数 $n(n+1)$ 都是偶数，这提示一个全称命题。证明只需将 $n$ 按奇偶分类：若 $n$ 偶，第一因子被 $2$ 整除；若 $n$ 奇，第二因子被 $2$ 整除。结论依赖“两个相邻整数恰有一个偶数”，并不依赖已经计算过多少个例子。

\begin{example}[计算很多项仍不能替代证明]
有人从若干计算猜测 $n^2+n+41$ 对所有非负整数 $n$ 都是素数。怎样检验这一猜测？
\end{example}
\begin{solve}
先尝试让前两项被 $41$ 整除。取 $n=41$，得到 $41^2+41+41=41\cdot43$，已经否定猜测；取 $n=40$ 还可得到 $40^2+40+41=41^2$。反例应同时给出符合范围的输入和失败原因，而不是仅说“继续算下去可能不成立”。
\end{solve}

书写证明时可以先在草稿中倒推目标，再在正文中按逻辑依赖正向排列。证明“若 $P$ 则 $Q$”时，每一步应来自 $P$、已知定理或前一步；若使用目标 $Q$ 参与推导，应明确它是在反证假设中出现，或说明所做变形均可逆。后文将经常使用这样的组织方式：先说明要构造什么，再说明构造为什么可行，最后验证结果。

\section{习题}

本章的共同出发点是整数的离散性：最小元支持归纳与递降，有限分类支持抽屉原理。做题时先辨认量词和对象范围，再选择证明方法；计算观察可以提出猜想，不能代替对所有对象的论证。

\begin{enumerate}
  \item 写出“对每个正整数 $n$，存在素数 $p>n$”的否定，说明否定命题需要怎样的证据。
  \item 用归纳法证明 $\sum_{k=1}^n k=n(n+1)/2$，指出归纳起点与归纳假设。
  \item 证明从 $1,2,\ldots,2n$ 中任取 $n+1$ 个不同整数，必有两个相邻。
  \item 仿照最小分母论证证明 $\sqrt3$ 无理；先说明平方被 $3$ 整除时，原整数也被 $3$ 整除。
 \item\label{ex:intro-negation} 写出“存在整数 $n$，使每个整数 $m$ 都满足 $m\le n$”的否定，并证明否定命题成立。
 \item\label{ex:intro-strong} 用强归纳法证明每个整数 $n\ge2$ 都能写成 $2a+3b$，其中 $a,b$ 为非负整数；说明为什么需要两个初始情形。
 \item\label{ex:intro-dividing-pair} 从 $1,2,\ldots,2n$ 中任取 $n+1$ 个不同整数，证明有两个数满足其中一个整除另一个。提示：不断除去因子 $2$。
 \item\label{ex:intro-descent-domain} 为什么正实数序列 $1,1/2,1/4,\ldots$ 可以无限下降，而正整数序列不能？指出递降法证明必须验证的范围条件。
\end{enumerate}

